% Figure 4. Molecules, fragments, complexes: the reaction plate. What the reader must see: one set of four nuclei
% grouped three ways along the reaction, KRb + KRb, the complex K2Rb2, K2 + Rb2, each grouping with its permutation
% group written by generators; the two product channels as two strands, gold for potassium and teal for rubidium,
% running side by side from the reactants into the complex and splitting at the products; the same two strands on
% the lattice of the sixteen subgroups of S, up from G_in to S and down to G_K and G_Rb, the complex above and the
% common core below. Colours: gold = potassium and its strand, teal = rubidium and its strand; ink elsewhere.
% Hatch unused. Positions schematic (declared); the lattice is computed (checks/verify.g).
\documentclass[10pt,border=1mm]{standalone}
\usepackage[T1]{fontenc}
\usepackage{figurestyle}
\input{primitives.tex}
\input{molecules.tex}
\input{numbers.tex}
\begin{document}
\begin{tikzpicture}[plate,kline/.style={line width=1.5pt,draw=ColGoldStroke},rbline/.style={line width=1.5pt,draw=ColTeal!85!PlateInk}]
\path[use as bounding box] (0,0) rectangle (15,11.0);
% ------------------------------------------------ A. the reaction: one set of four nuclei grouped three ways, the two strands
\foreach \x/\pname/\cap/\grp in {2.6/in/{$\mathrm{KRb}+\mathrm{KRb}$}/{$G_{\mathrm{in}}=\langle(12)(34),E^*\rangle$},7.5/complex/{$\mathrm K_2\mathrm{Rb}_2$}/{$S=\langle(12),(34),E^*\rangle$},12.4/out/{$\mathrm K_2+\mathrm{Rb}_2$}/{$G_{\mathrm K}=\langle(12),E^*\rangle$,\ $G_{\mathrm{Rb}}=\langle(34),E^*\rangle$}} {
  \begin{scope}[shift={(\x,9.0)}]
    \begin{scope}\molsolid\mollabelsinsideall\def\radK{.30}\def\radRb{.36}\pic[scale=.8] at (0,0) {mol3d-krb-\pname};\end{scope}
    \node[lab,anchor=north] at (0,-1.25) {\cap}; \node[sub,anchor=north] at (0,-1.68) {\grp};
  \end{scope}
}
\draw[kline] (4.5,9.08)--(5.5,9.08); \draw[rbline] (4.5,8.92)--(5.5,8.92);
\draw[kline] (9.5,9.08) .. controls (10.0,9.08) and (10.0,9.64) .. (10.5,9.64); \draw[rbline] (9.5,8.92) .. controls (10.0,8.92) and (10.0,8.36) .. (10.5,8.36);
% ------------------------------------------------ B. the sixteen subgroups of S with the same two strands
\begin{scope}[shift={(3.4,.9)}]
  \foreach \i/\x/\y/\m/\name in \krbNodes {\pgfmathsetmacro{\xx}{\x}\ifnum\i=10 \pgfmathsetmacro{\xx}{0.6429}\fi\ifnum\i=12 \pgfmathsetmacro{\xx}{0.3571}\fi \coordinate (k\i) at ({10.4*\xx},{1.65*\y});}
  \foreach \a/\b/\m in \krbCovers {\draw[line width=.3pt,draw=PlateInk!40] (k\a)--(k\b);}
  \draw[kline] ($(k\krbIdxIn)+(-.06,0)$)--($(k\krbIdxS)+(-.06,0)$); \draw[rbline] ($(k\krbIdxIn)+(.06,0)$)--($(k\krbIdxS)+(.06,0)$);
  \draw[kline] (k\krbIdxS)--(k\krbIdxK); \draw[rbline] (k\krbIdxS)--(k\krbIdxRb);
  \draw[line width=.7pt,draw=PlateInk] (k\krbIdxCore)--(k\krbIdxIn) (k\krbIdxCore)--(k\krbIdxK) (k\krbIdxCore)--(k\krbIdxRb) (k\krbIdxE)--(k\krbIdxCore);
  \foreach \i/\x/\y/\m/\name in \krbNodes {\ifnum\m=1\else\fill[PlateInk!45] (k\i) circle[radius=1.1pt];\fi}
  \draw[line width=.5pt,draw=PlateInk,fill=white] (k\krbIdxE) circle[radius=2.1pt]; \draw[line width=.5pt,draw=PlateInk,fill=white] (k\krbIdxCore) circle[radius=2.1pt]; \draw[line width=.5pt,draw=PlateInk,fill=white] (k\krbIdxS) circle[radius=2.6pt];
  \draw[line width=.5pt,draw=PlateInk,fill=ColPotassium] (k\krbIdxK) circle[radius=2.6pt]; \draw[line width=.5pt,draw=PlateInk,fill=ColRubidium] (k\krbIdxRb) circle[radius=2.6pt];
  \begin{scope}\clip (k\krbIdxIn) circle[radius=2.6pt]; \fill[ColPotassium] ($(k\krbIdxIn)+(-.2,-.2)$) rectangle ($(k\krbIdxIn)+(0,.2)$); \fill[ColRubidium] ($(k\krbIdxIn)+(0,-.2)$) rectangle ($(k\krbIdxIn)+(.2,.2)$);\end{scope} \draw[line width=.5pt,draw=PlateInk] (k\krbIdxIn) circle[radius=2.6pt];
  \node[lab,anchor=south] at ($(k\krbIdxS)+(0,.14)$) {$S$}; \node[lab,anchor=north] at ($(k\krbIdxE)+(0,-.14)$) {$E$};
  \node[lab,anchor=west] at ($(k\krbIdxCore)+(.14,0)$) {$\langle E^*\rangle$};
  \node[lab,anchor=north] at ($(k\krbIdxK)+(0,-.14)$) {$G_{\mathrm K}$}; \node[lab,anchor=north] at ($(k\krbIdxRb)+(0,-.14)$) {$G_{\mathrm{Rb}}$}; \node[lab,anchor=north] at ($(k\krbIdxIn)+(0,-.14)$) {$G_{\mathrm{in}}$};
  \foreach \y/\l in {0/{order $1$},1/{$2$},2/{$4$},3/{$8$}} \node[sub,anchor=east] at (-.6,{1.65*\y}) {\l};
\end{scope}
\end{tikzpicture}
\end{document}
